\documentclass[11pt]{article}
\usepackage[margin=1in]{geometry}
\usepackage{amsmath,amssymb}
\usepackage{graphicx}
\usepackage{booktabs}
\usepackage{xcolor}
\usepackage[colorlinks=true,linkcolor=blue!60!black,citecolor=blue!60!black,urlcolor=blue!60!black]{hyperref}

\title{Writing Together with Weave}
\author{Your Name \\ \small University of Chittagong}
\date{\today}

\begin{document}
\maketitle

\begin{abstract}
This is a starter article. Everyone invited to the project can edit it at the same time,
even offline; changes merge automatically when you reconnect.
\end{abstract}

\section{Introduction}
Weave keeps every collaborator's edits in a conflict-free replicated data type
(CRDT)~\cite{shapiro2011crdt}, so there is no central lock and no ``someone else is editing''
dialog.

\section{Some mathematics}
The probability of solving a problem of rating $r$ for a user of rating $R$ is
\begin{equation}
  P(\text{solve}) = \frac{1}{1 + 10^{(r - R)/400}}.
  \label{eq:elo}
\end{equation}
Equation~\eqref{eq:elo} is the Elo curve.

\section{A table}
\begin{table}[h]
  \centering
  \begin{tabular}{lrr}
    \toprule
    Algorithm & Interleaving & Complexity \\
    \midrule
    RGA  & backward only & $O(n)$ per op \\
    YATA & rare          & $O(n)$ per op \\
    Fugue & none         & $O(\log n)$ per op \\
    \bottomrule
  \end{tabular}
  \caption{Sequence CRDTs at a glance.}
\end{table}

\section{Conclusion}
Press \textbf{Ctrl/Cmd + S} to compile, and share the project from the top bar.

\bibliographystyle{plain}
\bibliography{refs}
\end{document}
